Este trabalho implementou, sem bibliotecas de controle, uma ferramenta de análise de estabilidade baseada no critério de Nyquist e a comparou com Routh-Hurwitz, Lugar das Raízes e Diagrama de Bode. O módulo de Nyquist trata polos sobre o eixo imaginário por indentação, avalia o atraso de transporte de forma exata e conta os envoltórios por dois algoritmos independentes que sempre concordaram. A ferramenta foi validada por 74 testes unitários, pela comparação com python-control e SciPy e por simulação temporal.

Nos quatro sistemas de teste, os métodos concordaram na fronteira de estabilidade até a nona casa decimal, o que indica que as diferenças práticas estão no que cada método exige e entrega, e não na resposta. O critério de Nyquist foi o único aplicado sem modificação a todos os casos. Com a planta instável em malha aberta, exigiu o envoltório anti-horário previsto pela relação $Z=N+P$. Com o atraso de rede, forneceu o atraso crítico exato de \TauCrit~s, enquanto a aproximação de Padé de primeira ordem errou esse valor em \ErroPadeApct\%, de forma não conservadora.

O estudo de desempenho mostrou que o custo do Nyquist é dominado pela amostragem do contorno e quase não depende da ordem, que o Routh é o mais leve e o Lugar das Raízes o mais caro. A precisão dos métodos baseados em raízes se degrada com a ordem por causa do mau condicionamento: o erro chega a $\ErroCompWiDoisZero$ para o polinômio de Wilkinson de grau 20. Um resultado relevante para a implementação é que um critério ingênuo de ``zero numérico'' no Routh produziu até \RobAbsErro\% de vereditos errados com aproximações de Padé de ordem alta, sem nenhum aviso; o escalonamento de frequência e uma tolerância relativa ao cancelamento eliminaram o problema.

Na aplicação de Engenharia da Computação, um laço PID com atraso de comunicação variável foi executado em dois Arduinos Uno no Tinkercad. O comportamento observado coincidiu amostra a amostra com o gêmeo digital, e o limite de estabilidade medido, entre \TkQuatroaD\ e \TkQuatrobD\ amostras, ficou a menos de uma amostra da previsão do critério de Nyquist. Com atraso variável e entrega em ordem, a estabilidade passou a ser governada pelo atraso mínimo, e não pelo médio.

Como trabalhos futuros, propõem-se: repetir a montagem em hardware físico, com atraso de rede real e perda de pacotes; projetar o controlador a partir da margem de atraso exigida; e substituir a grade fixa do cálculo de margens por um método adaptativo com garantia de detecção de todos os cruzamentos.
